Ce m\'emoire \'etudie la possibilit\'e de repr\'esenter un segment vid\'eo long par une structure textuelle suffisamment riche pour permettre \`a un grand mod\`ele de langage textuel d'identifier automatiquement un moment saillant. La contribution centrale du travail est un pipeline qui transforme un segment vid\'eo en repr\'esentation textuelle structur\'ee combinant transcription, segmentation en sc\`enes, association personnes-visages-locuteurs, indices de cadrage et r\'esum\'es visuels issus d'un mod\`ele vision-langage. Cette repr\'esentation est ensuite utilis\'ee comme contexte d'entr\'ee pour un mod\`ele s\'electionneur de clips.

L'\'evaluation repose sur un benchmark de 29 paires compos\'ees d'un segment source et d'un clip de r\'ef\'erence correspondant. Les r\'esultats consolid\'es actuellement montrent qu'une condition textuelle structur\'ee compl\`ete atteint un score combin\'e moyen de 0.301 sur l'ensemble complet, contre 0.273 pour une condition baseline en entr\'ee vid\'eo directe. Le m\'emoire adopte un format combin\'e centr\'e sur un article unique; les chapitres entourant l'article situent la contribution, documentent le benchmark et discutent le transfert vers une application de production.
